\chapter{Introduction}
\label{ch:introduction}

\section{Background and Motivation}
The rotary inverted pendulum, commonly called the Furuta pendulum, is an
underactuated mechanical system with one actuated rotary arm and one passive
pendulum joint. Its nonlinear, unstable, and coupled dynamics make it a useful
benchmark for modelling, state-feedback design, nonlinear control, sensing,
real-time implementation, and experimental validation
\cite{Furuta1991,Hamza2019}.

The experimental difficulty is not limited to deriving a stabilizing control
law. A real platform must estimate angular rates from noisy measurements,
handle wrapped encoder angles, apply commands through a finite-bandwidth
actuator, remain inside mechanical limits, and record enough information to
explain both successful and failed trials. These implementation effects are
especially important in the present rig because the arm is driven by a stepper
motor. The base angle is measured directly, but the motor still receives
open-loop STEP/DIR commands and can stall or miss motion under load.

This thesis therefore treats modelling, controller design, firmware, safety,
and experimental interpretation as one integrated research problem.

\section{Problem Statement}
The problem addressed in this thesis is to stabilize a custom stepper-driven
rotary inverted pendulum near the upright equilibrium while regulating or
commanding the rotary-arm position. The controller must operate on measured
arm and pendulum angles, generate commands compatible with the stepper
velocity interface, and remain usable despite sensor noise, actuator limits,
timing jitter, and mechanical travel limits.

The work is restricted to \emph{upright-only stabilization}. The pendulum is
manually placed near upright before engagement; no swing-up controller is
designed or evaluated.

\section{Research Questions}
The thesis is organized around three questions:
\begin{enumerate}
  \item How can a first-principles Furuta pendulum model be adapted to an
  acceleration-command representation that matches the stepper firmware
  interface without claiming ideal actuator realization?
  \item How do linear full-state feedback, a hybrid sliding-mode controller,
  and a four-state sliding-mode controller behave on the same physical
  platform under hold, commanded-motion, and manual-disturbance conditions?
  \item Which practical sensing, actuation, and safety mechanisms are required
  to make these controllers repeatable and interpretable on real hardware?
\end{enumerate}

\section{Research Objectives}
The objectives are:
\begin{enumerate}
  \item derive and numerically specialize a nonlinear and upright-linearized
  model of the custom rig;
  \item construct a measured-axis experimental platform and a real-time
  estimation, actuation, safety, logging, and analysis pipeline;
  \item implement one linear and two sliding-mode upright controllers using a
  common state, reference, and actuator interface;
  \item demonstrate commanded base-position motion while maintaining upright
  balance;
  \item compare the controllers using a fixed, auditable experimental dataset
  while keeping all conclusions within the limits of that evidence.
\end{enumerate}

\section{Contributions}
The principal contributions of this work are:
\begin{itemize}
  \item an acceleration-input formulation connected explicitly to the
  firmware path ``commanded acceleration $\rightarrow$ commanded speed
  $\rightarrow$ STEP/DIR actuation'';
  \item implementation of three upright controllers on the same custom Furuta
  platform with encoder feedback and stepper actuation;
  \item a supervised trapezoidal reference generator that permits commanded
  arm motion while the linear controller balances the pendulum;
  \item practical protections for wrapped angles, derivative noise, sensor
  glitches, actuator reversal, low-speed dithering, travel limits, and missed
  motion;
  \item a reproducible analysis path from pinned raw sessions to thesis
  figures, tables, and qualitative controller comparisons.
\end{itemize}

\section{Scope and Limitations}
The experimental evidence consists of seven pinned trials: three hold trials,
one linear-controller reference-tracking trial, and three manual tap trials.
Each controller--condition pair is represented by one trial. The dataset is
therefore suitable for a transparent engineering case study, but not for
statistical inference or a universal ranking of controller families.

The manual taps have no independent onset marker, so precise recovery-time or
impulse-response metrics cannot be claimed. The model also omits friction,
backlash, missed steps, motor-current dynamics, sensor faults, and saturation.
Those effects are handled only through firmware protections and experimental
interpretation.

\section{Thesis Organization}
Chapter~\ref{ch:literature} reviews Furuta modelling, linear and sliding-mode
control, experimental platforms, and the research gap addressed here.
Chapter~\ref{ch:platform_modelling} presents the hardware, notation, and
mathematical model. Chapter~\ref{ch:controller_implementation} develops the
three controllers and the real-time implementation. Chapter~\ref{ch:experiments}
defines the experimental and analysis methodology.
Chapter~\ref{ch:results} presents the measured results, and
Chapter~\ref{ch:conclusions} states the conclusions, limitations, and future
work.
